% Môn: Vật lý
% Lớp: 12
% Chương: Chương I. Vật lí nhiệt
% Bài: L12C1 Bài 5. Nhiệt nóng chảy riêng · Bài toán đun nóng chất lỏng đến nhiệt độ sôi rồi hóa hơi
% ===== Câu 21 =====
% ID: L12C1B5-62-DS
% Mức: NB
\dangbt{Bài toán đun nóng chất lỏng đến nhiệt độ sôi rồi hóa hơi}
\begin{ex}
% ID: L12C1B5-62-DS
% Mức: NB
Một chất lỏng ban đầu ở đúng nhiệt độ sôi.
\choiceTF
{\True $Q$ đun nóng đến sôi bằng 0.}
{\True Hóa hơi m_h cần Q=Lm_h.}
{\True Nếu hóa hơi toàn bộ thì m_h bằng khối lượng ban đầu.}
{Nhiệt độ tăng mạnh trong suốt quá trình sôi.}
\loigiai{
\begin{itemchoice}
\item (Đúng) $Q$ đun nóng đến sôi bằng 0.
\item (Đúng) Hóa hơi m_h cần Q=Lm_h.
\item (Đúng) Nếu hóa hơi toàn bộ thì m_h bằng khối lượng ban đầu.
\item (Sai) Nhiệt độ tăng mạnh trong suốt quá trình sôi.
\end{itemchoice}
}
\end{ex}

% ===== Câu 25 =====
% ID: L12C1B5-63-DS
% Mức: TH
\begin{ex}
% ID: L12C1B5-63-DS
% Mức: TH
Về hiệu suất trong quá trình đun và hóa hơi.
\choiceTF
{\True Nhiệt lượng hữu ích bằng mcΔt+Lm_h.}
{\True Nếu hiệu suất H<1 thì năng lượng nguồn lớn hơn nhiệt lượng hữu ích.}
{Có thể tính năng lượng nguồn bằng H.Q_hữu ích.}
{\True Hiệu suất phải đổi từ phần trăm sang số thập phân.}
\loigiai{
\begin{itemchoice}
\item (Đúng) Nhiệt lượng hữu ích bằng mcΔt+Lm_h.
\item (Đúng) Nếu hiệu suất H<1 thì năng lượng nguồn lớn hơn nhiệt lượng hữu ích.
\item (Sai) Có thể tính năng lượng nguồn bằng H.Q_hữu ích.
\item (Đúng) Hiệu suất phải đổi từ phần trăm sang số thập phân.
\end{itemchoice}
}
\end{ex}

% ===== Câu 26 =====
% ID: L12C1B5-74-TL
% Mức: NB
\begin{bt}
% ID: L12C1B5-74-TL
% Mức: NB
So sánh năng lượng để đun $1\,\text{kg}$ nước từ 20 ° $C$ đến 100 ° $C$ với năng lượng hóa hơi $1\,\text{kg}$ ở 100 °C.
\loigiai{
Q_đun=4200.80=3,36.10^ $5\,\text{J}$. Q_hóa hơi=2,3.10^ $6\,\text{J}$. Tỉ số khoảng 6,85; giai đoạn hóa hơi cần năng lượng lớn hơn nhiều.
}
\end{bt}

% ===== Câu 27 =====
% ID: L12C1B5-75-TLN
% Mức: NB
\begin{ex}
% ID: L12C1B5-75-TLN
% Mức: NB
Đun $0,50\,\text{kg}$ nước từ 40 ° $C$ đến 100 ° $C$ rồi hóa hơi $0,20\,\text{kg}$. Tính $Q$ theo kJ, lấy c=4200, L=2,3.10^6; chỉ ghi số.
\shortans{586}
\loigiai{
$Q1=0,5.4200.60=126 kJ; Q2=460 kJ; Q=586 kJ.$
}
\end{ex}

% ===== Câu 29 =====
% ID: L12C1B5-76-DS
% Mức: VD
\begin{ex}
% ID: L12C1B5-76-DS
% Mức: VD
Đun $1\,\text{kg}$ nước từ 20 ° $C$ đến 100 ° $C$ rồi hóa hơi $0,20\,\text{kg}$; c= $4200\,\text{J}$ /(kg.K), L=2,3.10^ $6\,\text{J}$ /kg.
\choiceTF
{\True Nhiệt lượng tăng nhiệt độ là 336 kJ.}
{\True Nhiệt lượng hóa hơi là 460 kJ.}
{\True Tổng nhiệt lượng là 796 kJ.}
{Có thể bỏ phần mcΔt.}
\loigiai{
\begin{itemchoice}
\item (Đúng) Nhiệt lượng tăng nhiệt độ là 336 kJ.
\item (Đúng) Nhiệt lượng hóa hơi là 460 kJ.
\item (Đúng) Tổng nhiệt lượng là 796 kJ.
\item (Sai) Có thể bỏ phần mcΔt.
\end{itemchoice}
}
\end{ex}

% ===== Câu 30 =====
% ID: L12C1B5-77-TL
% Mức: TH
\begin{bt}
% ID: L12C1B5-77-TL
% Mức: TH
Đun $2\,\text{kg}$ nước từ 25 ° $C$ đến 100 ° $C$ rồi hóa hơi $0,50\,\text{kg}$. Tính nhiệt lượng, lấy c= $4200\,\text{J}$ /(kg.K), L=2,3.10^ $6\,\text{J}$ /kg.
\loigiai{
Q1=2.4200.75=630000 $J; Q2=2,3.10^6.0,50=1150000\,\text{J}$. Tổng Q=1,78.10^ $6\,\text{J}$.
}
\end{bt}

% ===== Câu 31 =====
% ID: L12C1B5-78-TLN
% Mức: NB
\begin{ex}
% ID: L12C1B5-78-TLN
% Mức: NB
Tổng nhiệt lượng 630 kJ dùng đun $1,5\,\text{kg}$ chất có c=3000 từ 10 ° $C$ đến 70 ° $C$ rồi hóa hơi $0,40\,\text{kg}$. Tính $L$ theo đơn vị 10^ $5\,\text{J}$ /kg, chỉ ghi số.
\shortans{9}
\loigiai{
$Q1=270 kJ; Q_h=360 kJ; L=360000/0,40=9,0.10^5 J/kg.$
}
\end{ex}

% ===== Câu 33 =====
% ID: L12C1B5-79-DS
% Mức: VD
\begin{ex}
% ID: L12C1B5-79-DS
% Mức: VD
Đun $0,50\,\text{kg}$ nước từ 40 ° $C$ lên 100 ° $C$ rồi hóa hơi $0,10\,\text{kg}$.
\choiceTF
{\True Khối lượng trong phần mcΔt là $0,50\,\text{kg}$.}
{\True Khối lượng trong phần Lm là $0,10\,\text{kg}$.}
{Hai giai đoạn dùng cùng một công thức.}
{\True Nhiệt lượng tổng lớn hơn từng nhiệt lượng thành phần.}
\loigiai{
\begin{itemchoice}
\item (Đúng) Khối lượng trong phần mcΔt là $0,50\,\text{kg}$.
\item (Đúng) Khối lượng trong phần Lm là $0,10\,\text{kg}$.
\item (Sai) Hai giai đoạn dùng cùng một công thức.
\item (Đúng) Nhiệt lượng tổng lớn hơn từng nhiệt lượng thành phần. (Vì): ùng Q1=mcΔt và Q2=Lm; tổng Q=Q1+Q2
\end{itemchoice}
}
\end{ex}

% ===== Câu 34 =====
% ID: L12C1B5-80-TL
% Mức: TH
\begin{bt}
% ID: L12C1B5-80-TL
% Mức: TH
Một bài giải dùng Q=Lm cho nước từ 20 ° $C$ đến khi hóa hơi hết. Hãy chỉ ra sai sót và viết lời giải đúng dạng tổng quát.
\loigiai{
Sai vì bỏ qua nhiệt lượng tăng nhiệt độ từ 20 ° $C$ đến 100 °C. Đúng: Q=mc(100-20)+Lm=m[c.80+L], nếu bỏ qua nhiệt dung ấm và hao phí.
}
\end{bt}

% ===== Câu 35 =====
% ID: L12C1B5-81-TLN
% Mức: TH
\begin{ex}
% ID: L12C1B5-81-TLN
% Mức: TH
Một chất $2\,\text{kg}, c=2000\,\text{J}$ /(kg.K), tăng từ 25 ° $C$ đến 75 ° $C$ rồi hóa hơi $0,30\,\text{kg}$ với L=1,0.10^ $6\,\text{J}$ /kg. Tính $Q$ theo kJ, chỉ ghi số.
\shortans{500}
\loigiai{
$Q=2.2000.50+1,0.10^6.0,30=200+300=500 kJ.$
}
\end{ex}

% ===== Câu 37 =====
% ID: L12C1B5-82-DS
% Mức: VDC
\begin{ex}
% ID: L12C1B5-82-DS
% Mức: VDC
Xét phương pháp giải bài toán nhiều giai đoạn.
\choiceTF
{\True Cần chia quá trình theo từng giai đoạn.}
{\True Nhiệt lượng tổng bằng tổng đại số các phần.}
{Phải dùng khối lượng hóa hơi cho mọi giai đoạn.}
{\True Cần đổi đơn vị trước khi tính.}
\loigiai{
\begin{itemchoice}
\item (Đúng) Cần chia quá trình theo từng giai đoạn.
\item (Đúng) Nhiệt lượng tổng bằng tổng đại số các phần.
\item (Sai) Phải dùng khối lượng hóa hơi cho mọi giai đoạn.
\item (Đúng) Cần đổi đơn vị trước khi tính.
\end{itemchoice}
}
\end{ex}

% ===== Câu 38 =====
% ID: L12C1B5-83-TL
% Mức: VD
\begin{bt}
% ID: L12C1B5-83-TL
% Mức: VD
Một ấm chứa $1,2\,\text{kg}$ nước ở 30 °C. Sau khi nhận 1,2516 MJ hữu ích, nước đạt 100 ° $C$ và một phần hóa hơi. Tính khối lượng hóa hơi; c=4200, L=2,3.10^6.
\loigiai{
Q1=1,2.4200.70= $352800\,\text{J}. Q_h=1251600-352800=898800\,\text{J}. m_h=898800/(2,3.10^6)=0,3908\,\text{kg}$, xấp xỉ $0,391\,\text{kg}$.
}
\end{bt}

% ===== Câu 39 =====
% ID: L12C1B5-109-TLN
% Mức: VD
\begin{ex}
% ID: L12C1B5-109-TLN
% Mức: VD
Cấp 796 kJ để đun $1\,\text{kg}$ nước từ 20 ° $C$ đến 100 ° $C$ và hóa hơi một phần. c=4200, L=2,3.10^6. Khối lượng hóa hơi theo gam là bao nhiêu? Chỉ ghi số.
\shortans{18.7}
\loigiai{
$Q_h=796-336=460 kJ; m_h=460000/(2,3.10^6)=0,20 kg=200 g.$
}
\end{ex}

% ===== Câu 42 =====
% ID: L12C1B5-110-TL
% Mức: VD
\begin{bt}
% ID: L12C1B5-110-TL
% Mức: VD
Nêu cách xử lí nếu bài toán có thêm nhiệt dung của bình C_b và hiệu suất H.
\loigiai{
Nhiệt hữu ích gồm Q_h= $(mc+C_b)(t_s-t_0)+Lm_h. Năng lượng nguồn phải cung cấp là E=Q_h/$ H $, với$ H $viết dưới dạng số thập phân.$
}
\end{bt}

% ===== Câu 43 =====
% ID: L12C1B5-111-TLN
% Mức: VDC
\begin{ex}
% ID: L12C1B5-111-TLN
% Mức: VDC
Đun $1\,\text{kg}$ nước từ 20 ° $C$ đến 100 ° $C$ rồi hóa hơi $0,10\,\text{kg}$. c= $4200\,\text{J}$ /(kg.K), L=2,3.10^ $6\,\text{J}$ /kg. Tính $Q$ theo kJ, chỉ ghi số.
\shortans{566}
\loigiai{
$Q=1.4200.80+2,3.10^6.0,10=336+230=566 kJ.$
}
\end{ex}

% ===== Câu 46 =====
% ID: L12C1B5-112-TL
% Mức: VDC
\begin{bt}
% ID: L12C1B5-112-TL
% Mức: VDC
Thiết lập công thức tính nhiệt lượng khi đun khối lượng m chất lỏng từ t0 đến ts rồi hóa hơi khối lượng mh.
\loigiai{
Giai đoạn 1: Q1=mc(ts-t0). Giai đoạn 2: Q2=Lmh. Bỏ qua hao phí, Q=Q1+Q2=mc(ts-t0)+Lmh. Nếu hóa hơi toàn bộ thì mh=m.
}
\end{bt}

% ===== Câu 47 =====
% ID: L12C1B5-113-TLN
% Mức: VDC
\begin{ex}
% ID: L12C1B5-113-TLN
% Mức: VDC
Đun $0,80\,\text{kg}$ chất có c=2500 từ 20 ° $C$ đến 80 ° $C$ cần thêm 360 kJ để hóa hơi. Biết L=1,2.10^6. Khối lượng hóa hơi theo gam là bao nhiêu? Chỉ ghi số.
\shortans{300}
\loigiai{
$m_h=360000/(1,2.10^6)=0,30 kg=300 g.$
}
\end{ex}

% ===== Câu 48 =====
% ID: L12C1B5-114-TN
% Mức: NB
\begin{ex}
% ID: L12C1B5-114-TN
% Mức: NB
Đun $1\,\text{kg}$ chất lỏng có c= $4200\,\text{J}$ /(kg.K) từ 20 ° $C$ đến 100 ° $C$ rồi hóa hơi $0,2\,\text{kg}$, L=2,3.10^ $6\,\text{J}$ /kg. Bỏ qua hao phí. Tổng nhiệt lượng là
\choice
{79.6 kJ}
{\True 796 kJ}
{7960 kJ}
{796 $J$}
\loigiai{
Q=mcΔt+Lm=1.4200.80+2,3.10^6.0,2=796 kJ.
}
\end{ex}

% ===== Câu 52 =====
% ID: L12C1B5-115-TN
% Mức: NB
\begin{ex}
% ID: L12C1B5-115-TN
% Mức: NB
Đun $1,5\,\text{kg}$ chất lỏng có c= $3000\,\text{J}$ /(kg.K) từ 10 ° $C$ đến 70 ° $C$ rồi hóa hơi $0,4\,\text{kg}$, L=0,9.10^ $6\,\text{J}$ /kg. Bỏ qua hao phí. Tổng nhiệt lượng là
\choice
{63 kJ}
{\True 630 kJ}
{6300 kJ}
{630 $J$}
\loigiai{
Q=mcΔt+Lm=1,5.3000.60+0,9.10^6.0,4=630 kJ.
}
\end{ex}

% ===== Câu 56 =====
% ID: L12C1B5-116-TN
% Mức: NB
\begin{ex}
% ID: L12C1B5-116-TN
% Mức: NB
Nếu nước ban đầu đã ở 100 ° $C$ thì trong công thức tổng hợp
\choice
{\True phần mcΔt bằng 0}
{phần Lm bằng 0}
{cả hai phần bằng 0}
{phải dùng nhiệt nóng chảy}
\loigiai{
Δt=0 nên không có phần đun nóng; chỉ còn nhiệt lượng hóa hơi.
}
\end{ex}

% ===== Câu 60 =====
% ID: L12C1B5-117-TN
% Mức: TH
\begin{ex}
% ID: L12C1B5-117-TN
% Mức: TH
Đun $0,5\,\text{kg}$ chất lỏng có c= $4200\,\text{J}$ /(kg.K) từ 30 ° $C$ đến 100 ° $C$ rồi hóa hơi $0,1\,\text{kg}$, L=2,3.10^ $6\,\text{J}$ /kg. Bỏ qua hao phí. Tổng nhiệt lượng là
\choice
{26.2 kJ}
{\True 262 kJ}
{2620 kJ}
{262 $J$}
\loigiai{
Q=mcΔt+Lm=0,5.4200.70+2,3.10^6.0,1=262 kJ.
}
\end{ex}

% ===== Câu 64 =====
% ID: L12C1B5-118-TN
% Mức: TH
\begin{ex}
% ID: L12C1B5-118-TN
% Mức: TH
Khi đun chất lỏng từ dưới nhiệt độ sôi đến sôi rồi hóa hơi một phần, biểu thức đúng là
\choice
{Q=Lm toàn bộ}
{\True $Q=mc(t_s-t_0)+Lm_h$}
{$Q=mc(t_s-t_0)$}
{Q=mcΔt-Lm_h}
\loigiai{
Tổng nhiệt lượng gồm phần tăng nhiệt độ toàn bộ chất lỏng và phần hóa hơi m_h.
}
\end{ex}

% ===== Câu 68 =====
% ID: L12C1B5-119-TN
% Mức: TH
\begin{ex}
% ID: L12C1B5-119-TN
% Mức: TH
Đun cùng khối lượng nước từ 20 ° $C$ đến 100 ° $C$ rồi hóa hơi hoàn toàn. Thành phần thường lớn hơn là
\choice
{mcΔt}
{\True Lm}
{luôn bằng nhau}
{không xác định vì L=0}
\loigiai{
Với nước, cΔt≈3,36.10^ $5\,\text{J}$ /kg nhỏ hơn L≈2,3.10^ $6\,\text{J}$ /kg.
}
\end{ex}

% ===== Câu 72 =====
% ID: L12C1B5-120-TN
% Mức: VD
\begin{ex}
% ID: L12C1B5-120-TN
% Mức: VD
Đun $2\,\text{kg}$ chất lỏng có c= $2000\,\text{J}$ /(kg.K) từ 25 ° $C$ đến 75 ° $C$ rồi hóa hơi $0,5\,\text{kg}$, L=1.10^ $6\,\text{J}$ /kg. Bỏ qua hao phí. Tổng nhiệt lượng là
\choice
{70 kJ}
{\True 700 kJ}
{7000 kJ}
{700 $J$}
\loigiai{
Q=mcΔt+Lm=2.2000.50+1.10^6.0,5=700 kJ.
}
\end{ex}

% ===== Câu 76 =====
% ID: L12C1B5-121-TN
% Mức: VD
\begin{ex}
% ID: L12C1B5-121-TN
% Mức: VD
Trong bài toán đun đến sôi rồi hóa hơi, khối lượng dùng trong mcΔt là
\choice
{khối lượng phần hơi}
{\True khối lượng chất lỏng ban đầu được đun}
{luôn bằng $1\,\text{kg}$}
{khối lượng bình chứa}
\loigiai{
Toàn bộ chất lỏng ban đầu được tăng nhiệt độ đến điểm sôi.
}
\end{ex}

% ===== Câu 80 =====
% ID: L12C1B5-122-TN
% Mức: VDC
\begin{ex}
% ID: L12C1B5-122-TN
% Mức: VDC
Đun $0,8\,\text{kg}$ chất lỏng có c= $2500\,\text{J}$ /(kg.K) từ 20 ° $C$ đến 80 ° $C$ rồi hóa hơi $0,3\,\text{kg}$, L=1,2.10^ $6\,\text{J}$ /kg. Bỏ qua hao phí. Tổng nhiệt lượng là
\choice
{48 kJ}
{\True 480 kJ}
{4800 kJ}
{480 $J$}
\loigiai{
Q=mcΔt+Lm=0,8.2500.60+1,2.10^6.0,3=480 kJ.
}
\end{ex}

% ===== Câu 84 =====
% ID: L12C1B5-123-TN
% Mức: VDC
\begin{ex}
% ID: L12C1B5-123-TN
% Mức: VDC
Một kg nước từ 20 ° $C$ đến 100 ° $C$ cần 336 kJ; hóa hơi $0,10\,\text{kg}$ cần 230 kJ. Tổng là
\choice
{106 kJ}
{230 kJ}
{336 kJ}
{\True 566 kJ}
\loigiai{
$Q=336+230=566 kJ.$
}
\end{ex}
